% !TEX program = uplatex
% !TEX root = summer_2026_final_review_sheet.tex
\documentclass[12pt,a4paper]{jsarticle}

\usepackage{newtxtext,newtxmath}
\usepackage{array}
\usepackage{booktabs}
\usepackage{calc}
\usepackage{enumitem}
\usepackage{geometry}
\usepackage{xparse}
\usepackage{ulem}
\geometry{margin=15mm}

\setlength{\parindent}{0pt}
\setlength{\baselineskip}{1.25em}
\setlength{\extrarowheight}{0.35em}
\setlist[enumerate]{leftmargin=1.8em, itemsep=2.5em, topsep=0.45em, parsep=0.4em, partopsep=0pt}

\newlength{\fieldboxrule}
\newlength{\fieldboxsep}
\setlength{\fieldboxrule}{0.4pt}
\setlength{\fieldboxsep}{3pt}

\newcommand{\fieldbox}[2]{%
  \begingroup
  \setlength{\fboxsep}{\fieldboxsep}%
  \setlength{\fboxrule}{\fieldboxrule}%
  \fbox{%
    \makebox[\dimexpr#1-2\fieldboxsep-2\fieldboxrule\relax][l]{\rule{0pt}{#2}}%
  }%
  \endgroup
}

\newcommand{\fitblank}[1]{\fieldbox{#1zw}{1.85em}}
\newcommand{\choicecircle}{\makebox[1.3em][c]{\large ○}}
\newcommand{\namefield}[1][5.5cm]{\underline{\makebox[#1][l]{\rule{0pt}{1.85em}}}}
% 得点欄（下線）。第1引数=高さ（省略時 1.55em）、第2引数=下線の幅（省略時 1.35cm）
\NewDocumentCommand{\scorebox}{O{1.55em} O{1.35cm}}{%
  \makebox[\dimexpr#2+0.6cm\relax][c]{\underline{\rule{0pt}{#1}\hspace{#2}}\,/\,100}%
}
\NewDocumentCommand{\totalscorebox}{O{1.55em} O{1.35cm}}{%
  \makebox[\dimexpr#2+0.6cm\relax][c]{\underline{\rule{0pt}{#1}\hspace{#2}}\,/\,500}%
}
% 記入欄（枠線）。第1引数=幅（省略時 0.95\linewidth）、第2引数=高さ（省略時 1.85em）
\NewDocumentCommand{\sheetline}{O{0.95\linewidth} O{1.85em}}{%
  \fieldbox{#1}{#2}%
}
\newcommand{\scoretable}{%
  \renewcommand{\arraystretch}{1.55}%
  \begin{tabular}{C{1.85cm}C{1.85cm}C{1.85cm}C{1.85cm}C{1.85cm}C{2.55cm}}
  \toprule
  国語 & 数学 & 英語 & 理科 & 社会 & 合計 \\
  \midrule
  \scorebox & \scorebox & \scorebox & \scorebox & \scorebox & \totalscorebox \\
  \bottomrule
  \end{tabular}%
}
\newcolumntype{C}[1]{>{\centering\arraybackslash}p{#1}}

\pagestyle{empty}

\begin{document}
\input{sections/summary/最終振り返りシート_T.tex}
% % !TEX program = uplatex
% !TEX root = ../../summer_2026_final_review_sheet.tex
% S版（未作成）
% !TEX root = ../../summer_2026_final_review_sheet.tex

% --- 提出欄 ---
\begin{center}
    {\large \textbf{最終振り返りシート}}
  \end{center}

  \vspace{0.2em}
  \begin{flushright}
    {\small
    ITTO個別指導学院 札幌東校\\[0.35em]
    菅原\ 陸\hspace{1.5zw}{\textsf{Riku Sugawara}}\\[0.35em]
    08.2026
    }
  \end{flushright}

  \vspace{0.2em}
  \noindent
  \begin{tabular*}{\linewidth}{@{\extracolsep{\fill}} ll}
    \textbf{提出期限：8月8日（土）講義後}
    &
    \textbf{氏名：}\namefield[5.5cm]
  \end{tabular*}

  \vspace{1em}
  % 本文（1〜4：1ページ目）
  \begin{enumerate}
    \item 一斉夏期講習を受講し終えた自分を100点満点で評価してください。また、その理由を簡単に記入してください。
      \begin{center}
      \scorebox\\[1.5em]
      \sheetline[0.95\linewidth][4em]
      \end{center}

    \item 一斉夏期講習を担当した先生を100点満点で評価してください。また、その理由を簡単に記入してください。
      \begin{center}
      \scorebox\\[1.5em]
      \sheetline[0.95\linewidth][4em]
      \end{center}

    \item 一斉講習全体を通して（授業内容や授業の仕方、宿題など全て含めた上で）の評価を100点満点で評価してください。
    また、その理由を簡単に記入してください。
      \begin{center}
      \scorebox\\[1.5em]
      \sheetline[0.95\linewidth][4em]
      \end{center}

    \item 一斉講習を受講して良かった、成長できたことを記入してください。
      \begin{center}
      \sheetline[0.95\linewidth][8em]
      \end{center}

    \item 一斉講習に対しての改善点、アドバイスを記入してください。
      \begin{center}
      \sheetline[0.95\linewidth][3em]
      \end{center}

      \newpage

    \item 次回の一斉講習に対しての希望、要望を記入してください。
      \begin{center}
      \sheetline[0.95\linewidth][3em]
      \end{center}

    \item これから授業をしてくださる個別の先生方に知っておいてもらいたいことがあれば記入してください。
      \begin{center}
      \sheetline[0.95\linewidth][3em]
      \end{center}

    \item 次回の一斉講習があれば参加したいですか？また、その理由を簡単に記入してください。
      \begin{center}
        参加したい \qquad\qquad\qquad\qquad 参加したくない\\[1.5em]
      \sheetline[0.95\linewidth][2em]
      \end{center}

    \item 一斉講習を受講し終えた自分に一言記入してください。また、これからの意気込みを記入してください。
      \begin{center}
      \sheetline[0.95\linewidth][4em]
      \end{center}

    \item 菅原に伝えたいことを記入してください。
      \begin{center}
      \sheetline[0.95\linewidth][8em]
      \end{center}
    \end{enumerate}
    \vspace{1.5em}

    \begingroup
    \setlength{\parindent}{1zw}
    \setlength{\parskip}{0pt}

    \indent 一斉講習の受講、本当にお疲れ様でした。よく授業について来てくれました。ありがとう。

    講習を受ける前と今とでは何か変わりましたか？知識はもちろんのこと、受験生としての自覚を強く持ってもらうことも講習の意図の一つでした。

    もう受動的な学習はやめにしましょう。我々と一緒に能動的な学習をしていきましょう。そのスタートとして、一斉講習最後の宿題を課します。
    \[\textbf{夏休み中に一斉講習の内容を定着させること}\]

    具体的な方法をアドバイスしておきます。\uwave{同じ資料を使って、同じ問題を何度も解きなさい}。配布した資料は全てデータとして持っていますので、何度でも印刷して差し上げます。
    講習の資料は1から全て作成しました。したがって、私の意図が存分に反映されています。そのため、講習の資料を使って復習することが一番良いと思います。
    いろいろなテキスト、いろいろな問題をやって定着させたい気持ちは分かります。しかし、ここは私を信用してみてください。

    何をして良いか分からないときは \uwave{先生方を頼って} ください。必ず力になってくれます。

    今後の皆さんの成長を期待しています。

    \begin{flushright}
    菅原
    \end{flushright}

    \endgroup

\end{document}
